% !PRINTED 303
\weberpage{303}
%<<CONT
tained. But not all linear congruence groups are, conversely, also metacyclic, and to separate these off from them is yet a further problem, which up to now has been solved only in special cases $^{1)}$.

But in order at least to formulate the problem more precisely, we add the following considerations.

We denote the complete linear congruence group (6) by $R$, and separate off from it the group, contained within it, of the linear homogeneous substitutions

\weq{7}{\begin{aligned}
z_1' &\equiv \alpha_{1,1}z_1 + \alpha_{1,2}z_2 + \cdots + \alpha_{1,k}z_k \\
z_2' &\equiv \alpha_{2,1}z_1 + \alpha_{2,2}z_2 + \cdots + \alpha_{2,k}z_k \quad (\mathrm{mod}\ p) \\
&\cdots\cdots\cdots\cdots\cdots\cdots\cdots\cdots\cdots\cdots \\
z_k' &\equiv \alpha_{k,1}z_1 + \alpha_{k,2}z_2 + \cdots + \alpha_{k,k}z_k
\end{aligned}}

with non-vanishing determinant, which is called the \gesperrt{linear homogeneous congruence group}, and which we shall denote by $S$. Further, in (6) there is contained the $k$-fold cyclic group $Q$, which the substitutions

\weq{8}{z_1' = z_1 + \alpha_1, \quad z_2' = z_2 + \alpha_2, \dots, \quad z_k' = z_k + \alpha_k}

constitute.

If now $\sigma$ denotes an element of $S$ and $\gamma$ an element of $Q$, it follows without difficulty from the rule of composition \S.\,77, (7) that every substitution $\pi$ of the group $R$ can be represented \gesperrt{in one and only one way} in each of the two forms

\weq{9}{\pi = \sigma\gamma = \gamma'\sigma}

where, if $\sigma, \gamma$ are determined by (7), (8), $\gamma'$ denotes the substitution

\weq{10}{z_1' = z_1 + \alpha_1', \quad z_2' = z_2 + \alpha_2', \dots, \quad z_k' = z_k + \alpha_k'}

where the $\alpha_i'$ are derived from the $\alpha_i$ by the substitution $\sigma$.

$Q$ is a normal divisor of $R$; for if $\gamma_1$ is an arbitrary element of $Q$, then by repeated application of (9):

\[
\sigma\gamma\gamma_1 = \gamma'\gamma_1'\sigma = \gamma'\gamma_1'\sigma\gamma^{-1}\gamma = \gamma'\gamma_1'\gamma'^{-1}\sigma\gamma,
\]

hence, if

\[
\sigma\gamma = \pi, \quad \gamma'\gamma_1'\gamma'^{-1} = \gamma_1''
\]

is put,

\[
\pi\gamma_1 = \gamma_1''\pi,
\]
%CONT>>

\begin{center}
\rule{4cm}{0.4pt}
\end{center}

$^{1)}$ \gesperrt{C. Jordan}, Liouville's Journ. de Math\'em. S\'er. (2), Bd. XIII.
